\chapter{Jumps and Lévy Models}\label{ch:dv:jumps-and-levy-models}

A biotechnology company will publish the results of its main clinical trial on Thursday. Its one-week
at-the-money options trade at an implied volatility of 160\%. No one believes the share will diffuse at
160\% a year. It will drift at its ordinary 35\% until Thursday, then jump up about 22\% or down about 15\%,
and drift again. A diffusion needs one enormous volatility to reach the same price, and it then misprices
every other strike and every other expiry and gives the wrong hedge. A model with a jump needs the one
number the desk actually has an opinion on, the size of the move. This chapter covers the evidence that
index and single-stock options contain jumps and the three jump-diffusion models in use, Merton's, Kou's and
Bates's. It then turns to the pure-jump Lévy models, to what no hedge in the underlying can remove, and to the
situations in which jumps change the price of a trade.

\section{Why diffusions fail at short expiries}

In a diffusion, however volatile, the price moves continuously. Over a short time $T$ its log-return is
approximately normal with standard deviation $\sigma\sqrt T$, so an option struck a fixed distance out of the
money needs a move of many standard deviations to pay off. Its price falls faster than any power of $T$ as
the expiry shrinks: the normal tail $e^{-k^2/(2\sigma^2T)}$ decays exponentially in $1/T$. With jumps, a single
jump can carry the price past the strike at any time. The probability of one jump before $T$ is about
$\lambda T$ for a jump intensity $\lambda$, so the option's price is proportional to $T$ for small $T$:
\[
P_{\mathrm{OTM}}(T)\approx\lambda T\,\E\bigl[(K-S_0e^{J})^+\bigr]\quad\text{for a put, }T\to0,
\]
with $J$ the log-size of a jump. At the money both components contribute. The diffusion gives a price
proportional to $\sqrt T$, which dominates. Carr and Wu turned the difference in these speeds into a test of
what kind of process underlies the options, and found on S\&P~500 options both a continuous component and a
jump component.

\begin{example}[A ten-percent put as expiry shrinks]\label{ex:dv:jumps-and-levy-models:carrwu}
Take the Merton model fitted below to chapter 9's surface ($\sigma=9.8\%$, $\lambda=3.16$ jumps a year of mean
log-size $-5.8\%$). Its ten-percent out-of-the-money put, divided by its expiry, tends to
$\lambda\E[(K-S_0e^J)^+]=1.00$ a year: at one day the ratio is 1.02. Under Black--Scholes with the same total
variance, 16.4\% a year, the same put is worth $3\times10^{-36}$ at one day and $6\times10^{-9}$ at five days
(\cref{fig:dv:jumps-and-levy-models:short}, left).
\end{example}

The same argument explains the smile. A diffusion's short-dated implied volatility is flat, because its
short-dated distribution is normal. A jump's is steep, because its short-dated distribution is a mixture of
``nothing happened'' and ``one jump happened'', with fat tails. For a Lévy process, the class that contains
every model of this chapter except Bates's, the cumulants of the log-return all grow linearly in $T$. The
skewness therefore scales like $T^{-1/2}$ and the excess kurtosis like $T^{-1}$. The smile is steepest at short
expiries and fades as the central limit theorem takes over.

\begin{omfigure}
\begin{tikzpicture}
\begin{groupplot}[group style={group size=2 by 1,horizontal sep=1.6cm},
  width=6.6cm,height=5cm,
  tick label style={font=\small},label style={font=\small},grid=major,grid style={gray!25},table/col sep=comma]
\nextgroupplot[xmode=log,ymode=log,xlabel={expiry $T$ (years)},ylabel={price of the 90 put},
  xmin=0.0025,xmax=0.28,ymin=1e-8,ymax=3,
  title={\small a ten-percent put},
  legend columns=1,legend cell align=left,legend style={font=\footnotesize,at={(0.5,-0.32)},anchor=north,draw=none,fill=none}]
\addplot[omThm,very thick] table[x=t,y=jump]{figdata/derivatives/13-jumps-and-levy-models/otm_decay.csv};
\addplot[omDef,very thick,dashed,unbounded coords=jump] table[x=t,y=bs]{figdata/derivatives/13-jumps-and-levy-models/otm_decay.csv};
\addplot[gray,thick,dotted] table[x=t,y=limit]{figdata/derivatives/13-jumps-and-levy-models/otm_decay.csv};
\legend{Merton,{Black--Scholes, same variance},{$\lambda T\,\E[(K-S_0e^J)^+]$}}
\nextgroupplot[xmode=log,ymode=log,xlabel={expiry $T$ (years)},ylabel={at-the-money skew $|\psi(T)|$},
  xmin=0.015,xmax=2.6,ymin=0.025,ymax=5,
  xtick={0.02,0.1,0.5,2},xticklabels={0.02,0.1,0.5,2},ytick={0.03,0.1,0.3,1,3},yticklabels={0.03,0.1,0.3,1,3},
  title={\small skew term structure},
  legend columns=2,legend cell align=left,legend style={font=\footnotesize,at={(0.5,-0.32)},anchor=north,draw=none,fill=none}]
\addplot[black,only marks,mark=*,mark size=1.8pt] table[x=t,y=market]{figdata/derivatives/13-jumps-and-levy-models/skews.csv};
\addplot[omThm,very thick] table[x=t,y=merton]{figdata/derivatives/13-jumps-and-levy-models/skews.csv};
\addplot[omDef,thick,dashed] table[x=t,y=vg]{figdata/derivatives/13-jumps-and-levy-models/skews.csv};
\addplot[omProp,thick] table[x=t,y=heston]{figdata/derivatives/13-jumps-and-levy-models/skews.csv};
\addplot[gray,very thick,dotted] table[x=t,y=bates]{figdata/derivatives/13-jumps-and-levy-models/skews.csv};
\legend{market,Merton,variance gamma,Heston,Bates}
\end{groupplot}
\end{tikzpicture}

\omcaption{Short expiries separate jumps from diffusion. Left: a ten-percent out-of-the-money put under the
fitted Merton model is proportional to its expiry; under Black--Scholes with the same variance it vanishes
faster than any power. Right: the at-the-money skew of chapter 9's surface against Merton and variance gamma
fitted at one month, Heston fitted to the surface (chapter 10) and Bates fitted to all eight expiries. Data:
the tutorial.}\label{fig:dv:jumps-and-levy-models:short}
\end{omfigure}

\section{Jump-diffusions: Merton, Kou, Bates}

A jump-diffusion adds to the Brownian motion of chapter 3 a compound Poisson process (One Quant Book 4,
chapter 6): jumps arrive at rate $\lambda$ and each multiplies the price by $e^{J}$, with the $J$ independent
and identically distributed. To keep the discounted price a martingale, the drift is lowered by the expected
relative jump, $\lambda\bar k$ with $\bar k=\E[e^J]-1$. With $x_T=\ln(S_T/F_T)$ and $\phi_J$ the
characteristic function of $J$, every such model has
\[
\E\bigl[e^{iux_T}\bigr]=\exp\Bigl(T\bigl(-\tfrac12\sigma^2(iu+u^2)-iu\lambda\bar k+\lambda(\phi_J(u)-1)\bigr)\Bigr),
\]
the Lévy--Khintchine form (Book 4, chapter 6) with a Gaussian part and a finite jump measure. The three
models differ in the law of $J$, and in whether the volatility is itself random.

\begin{definition}[Merton jump-diffusion model]\label{def:dv:jumps-and-levy-models:merton}
The \emph{Merton jump-diffusion model}\index{Merton jump-diffusion model} is the jump-diffusion with normal
log-jumps, $J\sim\mathcal N(\mu_J,\delta^2)$, so $\phi_J(u)=e^{iu\mu_J-\frac12\delta^2u^2}$ and
$\bar k=e^{\mu_J+\frac12\delta^2}-1$. Its parameters are $(\sigma,\lambda,\mu_J,\delta)$.
\end{definition}

Merton's model has a closed-form price. Condition on the number $n$ of jumps before $T$, which is Poisson with
mean $\lambda T$. Given $n$, the log-return is normal, so the option is a Black price with a shifted forward and
a larger variance:
\[
C=\sum_{n\ge0}e^{-\lambda T}\frac{(\lambda T)^n}{n!}\,C_{\mathrm{Black}}\Bigl(F\,e^{n(\mu_J+\frac12\delta^2)-\lambda\bar kT},\,K,\,T,\,
\sqrt{\sigma^2+n\delta^2/T}\Bigr).
\]
The tutorial checks it against the Fourier price of chapter 10's integral to a few millionths.

\begin{definition}[Kou model]\label{def:dv:jumps-and-levy-models:kou}
The \emph{Kou model}\index{Kou model} is the jump-diffusion with double-exponential log-jumps: with probability
$p$ the jump is up and exponential with rate $\eta_1>1$, otherwise down and exponential with rate $\eta_2$, so
$\phi_J(u)=p\,\frac{\eta_1}{\eta_1-iu}+(1-p)\frac{\eta_2}{\eta_2+iu}$. Its parameters are
$(\sigma,\lambda,p,\eta_1,\eta_2)$.
\end{definition}

Kou chose the double exponential for two reasons. It gives asymmetric tails that are heavier than normal on
each side, and the exponential law is memoryless: the overshoot of a jump over a barrier is again exponential,
which yields analytical prices for path-dependent options as well as vanillas. The condition $\eta_1>1$ keeps $\E[e^J]$ finite.

\begin{definition}[Bates model]\label{def:dv:jumps-and-levy-models:bates}
The \emph{Bates model}\index{Bates model} is Heston's \omterm{def:dv:stochastic-volatility:sv}{stochastic volatility model} (chapter 10) with Merton jumps
added to the price. The characteristic function is the product of Heston's and the jump factor, and the
parameters are $(v_0,\kappa,\bar v,\eta,\rho,\lambda,\mu_J,\delta)$.
\end{definition}

The division of labour is the point of Bates's model. Jumps make the short-dated skew. Stochastic
volatility makes the long-dated skew and the term structure. Neither can do the other's job: jump skew fades
like $T^{-1/2}$, and Heston's short skew levels off (chapter 10).

\begin{example}[Three models on one expiry, one on all]\label{ex:dv:jumps-and-levy-models:fits}
Fitted to the one-month smile of chapter 9's surface (nine strikes within two standard moves), Merton returns
$\sigma=9.8\%$, $\lambda=3.16$, $\mu_J=-5.8\%$, $\delta=4.6\%$ with a root-mean-square error of 0.18 volatility
point. Kou returns $\sigma=8.4\%$, $\lambda=11.7$, $p=0.24$, $\eta_1=85$, $\eta_2=30.4$, with an error of 0.06.
Variance gamma (next section) has an error of 0.24 (\cref{fig:dv:jumps-and-levy-models:fits}). Carried to other
expiries, all three fail. At one year Merton's skew is $-0.07$ against the market's $-0.25$. At one week variance
gamma's is $-3.66$ against $-1.44$ (\cref{fig:dv:jumps-and-levy-models:short}). Bates fitted to all eight expiries
from one week to two years has an error of 0.34 point, with $v_0=0.018$, $\kappa=2.06$, $\bar v=0.062$,
$\eta=0.85$, $\rho=-0.70$ and jumps $\lambda=2.34$, $\mu_J=-3.6\%$, $\delta=3.0\%$. Its skew is $-1.36$ at one week
and $-0.23$ at one year.
\end{example}

\begin{omfigure}
\begin{tikzpicture}
\begin{axis}[width=10.5cm,height=5cm,
  xlabel={strike},ylabel={1-month implied volatility (\%)},
  tick label style={font=\small},label style={font=\small},
  xmin=87,xmax=114,ymin=10,ymax=24,grid=major,grid style={gray!25},
  legend columns=3,legend cell align=left,legend style={font=\footnotesize,at={(0.5,-0.3)},anchor=north,draw=none,fill=none,/tikz/every even column/.append style={column sep=8pt}},
  table/col sep=comma]
\addplot[black,only marks,mark=*,mark size=1.8pt] table[x=k,y=market]{figdata/derivatives/13-jumps-and-levy-models/fit_market.csv};
\addplot[omThm,very thick] table[x=k,y=merton]{figdata/derivatives/13-jumps-and-levy-models/fit_models.csv};
\addplot[omDef,thick,dashed] table[x=k,y=kou]{figdata/derivatives/13-jumps-and-levy-models/fit_models.csv};
\addplot[omProp,thick,dotted] table[x=k,y=vg]{figdata/derivatives/13-jumps-and-levy-models/fit_models.csv};
\addplot[gray,thick] table[x=k,y=heston]{figdata/derivatives/13-jumps-and-levy-models/fit_models.csv};
\legend{market quotes,Merton,Kou,variance gamma,Heston (chapter 10)}
\end{axis}
\end{tikzpicture}

\omcaption{One-month smiles of chapter 9's surface: the nine fitted quotes, the three jump models fitted to
them, and Heston calibrated to the whole surface. On one expiry every jump model fits; the differences appear
at other expiries. Data: the tutorial.}\label{fig:dv:jumps-and-levy-models:fits}
\end{omfigure}

\section{Pure-jump Lévy models}

A Lévy process with infinitely many small jumps in every interval can do without the Brownian part. The small
jumps play the role of the diffusion, and the large ones make the tails. The best known such model changes the
clock of a Brownian motion.

\begin{definition}[Variance gamma model]\label{def:dv:jumps-and-levy-models:vg}
The \emph{variance gamma model}\index{variance gamma model} runs a Brownian motion with drift $\theta$ and
volatility $\sigma$ on a random clock $g_t$, a gamma process with mean rate one and variance rate $\nu$ (a
subordinator, One Quant Book 4, chapter 6): $x_t=\omega t+\theta g_t+\sigma W_{g_t}$, with $\omega$ set so that
$\E[e^{x_t}]=1$. Its characteristic function is
$\E[e^{iux_t}]=e^{iu\omega t}\bigl(1-iu\theta\nu+\frac12\sigma^2\nu u^2\bigr)^{-t/\nu}$, and
$\omega=\frac1\nu\ln\bigl(1-\theta\nu-\frac12\sigma^2\nu\bigr)$.
\end{definition}

The three parameters have direct meanings. $\sigma$ sets the scale, $\nu$ the kurtosis (the variability of the
clock), and $\theta$ the skewness (the drift that the clock carries). Madan, Carr and Chang found that on
S\&P~500 data the statistical density was nearly symmetric while the \omterm{def:dv:implied-volatility-and-its-surface:density}{risk-neutral density} was negatively skewed
with more kurtosis. The skew of index options, on this reading, is a price of downside risk more than a
feature of realised returns. Conditioning on the clock gives a closed form. Given $g_T=g$ the log-return is
normal, so the call is a Black price averaged over the gamma law of $g$, which the build computes by quadrature
and checks against the Fourier price. The CGMY model of Carr, Geman, Madan and Yor generalises variance gamma
with a fourth parameter that controls how active the small jumps are. It covers finite and infinite activity,
and finite and infinite variation.

\begin{example}[Moments that fade]\label{ex:dv:jumps-and-levy-models:moments}
For the Merton model fitted above, the skewness of the log-return is $-2.90$ at one week, $-1.39$ at one month
and $-0.40$ at one year, and the excess kurtosis 15.2, 3.5 and 0.29: exactly the $T^{-1/2}$ and $T^{-1}$ laws. The
fitted \omterm{def:dv:jumps-and-levy-models:vg}{variance gamma model} is within a few percent of the same numbers, since both were fitted to the same
month. That is why a Lévy model fitted at one expiry cannot fit the others. The market's skew decays like
$T^{-0.44}$ (chapter 12), close to the Lévy exponent $-\frac12$ at short expiries, but at long expiries it is held
up by the persistence of volatility, which no Lévy process has.
\end{example}

\section{What cannot be hedged}

A market with jumps of random size is incomplete (chapter 1): the share alone cannot replicate an option.
Delta hedging works in a diffusion because over a short interval the option's value is linear in the move to
first order, and the second-order term averages out through gamma and \omterm{def:dv:greeks-and-the-hedging-pnl:greeks}{theta}. A jump is not small. Over a jump
the hedged position gains or loses the option's convexity on the whole move,
$V(Se^J)-V(S)-\Delta S(e^J-1)$, and no choice of $\Delta$ removes it for every $J$ at once.

\begin{example}[Delta hedging in a world with jumps]\label{ex:dv:jumps-and-levy-models:hedge}
Sell a one-month at-the-money call at the fitted Merton price, 1.89 on a spot of 100, and delta-hedge it daily
with Merton's own delta. Over 20\,000 paths the hedged P\&L has mean zero and a standard deviation of 0.97 times
the premium, and one path in a hundred loses more than 4.1 premiums. In a diffusion with the same total variance,
hedged the same way, the standard deviation is 0.185 premiums and the one-in-a-hundred loss 0.50
(\cref{fig:dv:jumps-and-levy-models:hedge}, left). The daily hedge is as good in both worlds between jumps; the
difference is the handful of days with a jump.
\end{example}

The scheduled event is the limiting case, and a revealing one. If the outcome has two known sizes, the event is
a one-step binomial tree (chapter 2): two states and one hedge instrument, a \omterm{def:dv:no-arbitrage-and-the-fundamental-theorems:complete}{complete market}. The ratio
$\Delta^*=\frac{V_u-V_d}{S_u-S_d}$ removes the jump risk entirely. It is not the option's delta, which measures
the response to small moves. Once the sizes are uncertain the tree has more branches than instruments, and a
residual remains that only options can hedge.

\begin{example}[Hedging across the trial result]\label{ex:dv:jumps-and-levy-models:event}
The desk is short the biotechnology share's one-week at-the-money straddle, worth 8.86 on a spot of 50 (a
diffusion at 35\%, an up move of 22.1\% with probability 0.4, a down move of 14.8\%). It hedges from just before
the announcement to just after. The standard deviation of the P\&L across the event is 1.81, or 20.4\% of the
premium, with no hedge. It is 2.01 with the model's own delta, $-0.02$, which responds only to small moves. It is
1.01 (11.4\%) with the Black delta at the straddle's implied volatility, $+0.09$. With the binomial ratio,
$+0.20$, it is zero. When each outcome's size is itself uncertain (a standard deviation of 4\% in log around each),
the best single ratio is still 0.20, but it leaves 1.45, 16.5\% of the premium
(\cref{fig:dv:jumps-and-levy-models:hedge}, right).
\end{example}

\begin{omfigure}
\begin{tikzpicture}
\begin{groupplot}[group style={group size=2 by 1,horizontal sep=1.6cm},
  width=6.6cm,height=5cm,
  tick label style={font=\small},label style={font=\small},grid=major,grid style={gray!25},table/col sep=comma]
\nextgroupplot[xlabel={hedged P\&L (premiums)},ylabel={share of paths},ymode=log,xmin=-5,xmax=1.5,ymin=1e-4,ymax=1,
  title={\small daily delta hedge, one month},
  legend columns=2,legend style={font=\footnotesize,at={(0.5,-0.32)},anchor=north,draw=none,fill=none}]
\addplot[omThm,very thick,const plot] table[x=x,y=merton]{figdata/derivatives/13-jumps-and-levy-models/hedge_hist.csv};
\addplot[omDef,thick,dashed,const plot] table[x=x,y=diffusion]{figdata/derivatives/13-jumps-and-levy-models/hedge_hist.csv};
\legend{Merton,diffusion}
\nextgroupplot[xlabel={hedge},ylabel={P\&L standard deviation (\% of premium)},ymin=0,ymax=25,xmin=-0.6,xmax=4.6,
  xtick=data,xticklabels={none,model,Black,binomial,uncertain},
  x tick label style={font=\scriptsize,rotate=35,anchor=north east},
  title={\small across the announcement}]
\addplot[omThm,ybar,bar width=12pt,fill=omThm!40] table[x=i,y=sd]{figdata/derivatives/13-jumps-and-levy-models/event_hedge.csv};
\end{groupplot}
\end{tikzpicture}

\omcaption{What a hedge in the share leaves. Left: the P\&L of a short one-month call hedged daily, in the fitted
Merton world and in a diffusion of equal variance (log scale: the jump world's left tail is the difference).
Right: the standard deviation of the P\&L of a short straddle hedged across a binary announcement, with no hedge,
the model's delta, the Black delta and the binomial ratio (known outcome sizes), and the best ratio when the sizes
are uncertain; the binomial ratio removes it all when the sizes are known. Data: the
tutorial.}\label{fig:dv:jumps-and-levy-models:hedge}
\end{omfigure}

\section{When jumps matter}

Jumps change prices and hedges wherever the payoff is sensitive to large moves over short times.
\begin{itemize}
\item \textbf{Scheduled events.} Earnings, trial results, elections, central-bank decisions: the \omterm{def:dv:parametrising-the-surface:event}{event variance}
of chapter 8 is a jump. A two-point model reads the probability and sizes from the smile, which a single implied
volatility cannot.
\item \textbf{Short-dated out-of-the-money options.} Their value is mostly the probability of a jump
(\cref{ex:dv:jumps-and-levy-models:carrwu}); a diffusion calibrated elsewhere makes them nearly free.
\item \textbf{Gap-sensitive payoffs.} A barrier monitored continuously (chapter 15) can be crossed by a jump
without the price trading near it, so the hedge the desk plans at the barrier never happens. Digitals and
autocallables near their triggers, and leveraged products that must be rebalanced, carry the same gap risk.
\item \textbf{Hedging reserves.} Where the hedge leaves a jump residual, the price must carry a reserve for it
(chapter 27); the size of the residual in \cref{ex:dv:jumps-and-levy-models:hedge} is the order of magnitude.
\end{itemize}
Where they matter less: long-dated vanillas, whose smile is mostly the dynamics of volatility, and at-the-money
options, whose price is mostly the diffusion.

\begin{omfigure}
\begin{tikzpicture}
\begin{groupplot}[group style={group size=2 by 1,horizontal sep=1.6cm},
  width=6.6cm,height=5cm,
  tick label style={font=\small},label style={font=\small},grid=major,grid style={gray!25},table/col sep=comma]
\nextgroupplot[xlabel={share price at expiry},ylabel={density},xmin=25,xmax=80,ymin=0,
  yticklabel style={/pgf/number format/fixed},
  title={\small one-week density},
  legend columns=1,legend cell align=left,legend style={font=\footnotesize,at={(0.5,-0.32)},anchor=north,draw=none,fill=none}]
\addplot[omThm,very thick] table[x=s,y=event]{figdata/derivatives/13-jumps-and-levy-models/event_density.csv};
\addplot[omDef,thick,dashed] table[x=s,y=lognormal]{figdata/derivatives/13-jumps-and-levy-models/event_density.csv};
\legend{diffusion and binary event,{lognormal at the event's 160\%}}
\nextgroupplot[xlabel={strike},ylabel={1-week implied volatility (\%)},xmin=37,xmax=63,ymin=60,ymax=170,
  title={\small one-week smile}]
\addplot[omThm,very thick,mark=*,mark size=1.5pt] table[x=k,y=vol]{figdata/derivatives/13-jumps-and-levy-models/event_smile.csv};
\end{groupplot}
\end{tikzpicture}

\omcaption{A binary announcement inside a one-week option on a share at 50: a 35\% diffusion plus a move of
$+22.1\%$ (probability 0.4) or $-14.8\%$. Left: the density has two humps, which no single lognormal
reproduces. Right: the implied volatility is 160\% at the money and falls away on both sides, the ``frown'' of a
bimodal distribution. Data: the chapter's code.}\label{fig:dv:jumps-and-levy-models:event}
\end{omfigure}

\begin{example}[Reading the event from the smile]\label{ex:dv:jumps-and-levy-models:read}
The at-the-money volatility of the biotechnology share, 160.4\%, read with chapter 8's rule against a 35\% diffusion,
gives an \omterm{def:dv:parametrising-the-surface:event}{event variance} of 0.047. The rule's standard deviation for the move is then 21.7\%, and its expected
absolute move is 17.3\%. The true move has a standard deviation of 17.6\% and an expected absolute size of 17.6\%.
The straddle prices the expected absolute move, which the rule gets right, while the rule's standard deviation
assumes a normal move and overstates the binary one. A two-point model fitted to the whole smile
(\cref{fig:dv:jumps-and-levy-models:event}) recovers the probability and both sizes. It does so even when the
market's outcome sizes are uncertain, as the weekend problem shows.
\end{example}

\section{Tutorial: one interface for jump models}\label{tut:dv:jumps-and-levy-models:tut}

\begin{tutorial}
\textbf{Goal.} Implement four jump models and the event model behind one characteristic-function interface,
check each against a closed form, fit them to a smile, and measure the hedging residuals. \textbf{End state:} the
four figures of the chapter and the numbers of the weekend problem.
\begin{tutsteps}
\item \textbf{A model is its characteristic function.} Merton's, martingale-corrected, with the series price
that checks it:
\omcode{code/firm/jumps/firm_jumps.py}{38}{67}{Merton's model: characteristic function, cumulants, series price.}\label{lst:dv:jumps-and-levy-models:merton}
\item \textbf{One pricer for all.} Lewis's integral from chapter 10, on a grid long enough for pure-jump models
at short expiries, whose characteristic functions decay only like a power of $u$:
\omcode{code/firm/jumps/firm_jumps.py}{175}{182}{Fourier pricer consuming any model's characteristic function.}\label{lst:dv:jumps-and-levy-models:lewis}
\item \textbf{The scheduled event} as a model with an exact two-point mixture price:
\omcode{code/firm/jumps/firm_jumps.py}{147}{172}{The binary event model.}\label{lst:dv:jumps-and-levy-models:event}
\item \textbf{Run} \texttt{dv\_jumps.fit\_models()}, \texttt{fit\_bates()}, \texttt{event\_hedge()},
\texttt{hedge\_experiment()} and \texttt{fig\_jumps.py}.
\end{tutsteps}
\textbf{What to change next.} Fit Merton at one week instead of one month and compare the one-year skew; set the
outcome-size uncertainty to 8\% and re-run the event hedge; replace the gamma clock of variance gamma by an
inverse Gaussian one (the normal inverse Gaussian model) and compare the fits.
\end{tutorial}

\section{Build: jump models behind one interface}\label{bld:dv:jumps-and-levy-models:jumps}

\begin{build}
\bfield{Purpose} The miniature firm's jump models, as characteristic functions that the transform engine of
chapter 24 prices and calibrates, and the scheduled-event model the options desk uses around announcements.
\bfield{Interface} Dataclasses \texttt{Merton}, \texttt{Kou}, \texttt{Bates}, \texttt{VarianceGamma},
\texttt{EventJump} (with \texttt{from\_up}), each with \texttt{cf(u, t)} of $\ln(S_t/F_t)$ and, where they exist,
\texttt{cumulants(t)}, \texttt{series\_call}, \texttt{conditional\_call}, \texttt{mixture\_call};
\texttt{call\_prices(model, fwd, strikes, t)}; \texttt{implied\_vols}; \texttt{skew\_kurtosis};
\texttt{calibrate(make, z0, quotes, fwd) -> (model, z, rmse)}.
\bfield{Rules} Every characteristic function satisfies $\phi(0)=1$ and $\phi(-i)=1$ (the martingale condition),
checked in the tests; parameters are fitted through transforms that keep them admissible; the Fourier grid is
checked against a closed form at the shortest expiry the desk trades.
\bfield{Acceptance tests} \texttt{code/firm/jumps/tests/}: the martingale condition for every model at one week
and one year; Fourier against Merton's series, variance gamma's conditional price and the event mixture; no
jumps gives Black--Scholes and Bates without jumps gives Heston; skewness and kurtosis scale as $T^{-1/2}$ and
$T^{-1}$; calibration recovers known Merton parameters.
\bfield{Stretch} CGMY and normal inverse Gaussian models with the same interface; Kou's closed-form barrier
prices as a check on chapter 15's engines.
\end{build}

\begin{omsources}
\item R.~C.~Merton, ``Option pricing when underlying stock returns are discontinuous'', \emph{Journal of Financial
Economics} 3(1--2) (1976) 125--144.
\item S.~G.~Kou, ``A jump-diffusion model for option pricing'', \emph{Management Science} 48(8) (2002) 1086--1101.
\item D.~S.~Bates, ``Jumps and stochastic volatility: exchange rate processes implicit in Deutsche mark options'',
\emph{Review of Financial Studies} 9(1) (1996) 69--107.
\item D.~B.~Madan, P.~P.~Carr and E.~C.~Chang, ``The variance gamma process and option pricing'', \emph{Review of
Finance} 2(1) (1998) 79--105.
\item P.~Carr, H.~Geman, D.~B.~Madan and M.~Yor, ``The fine structure of asset returns: an empirical
investigation'', \emph{Journal of Business} 75(2) (2002) 305--333.
\item P.~Carr and L.~Wu, ``What type of process underlies options? A simple robust test'', \emph{Journal of
Finance} 58(6) (2003) 2581--2610.
\end{omsources}

\section{Exercises}

\begin{exercise}[$\star$]\label{exo:dv:jumps-and-levy-models:1}
In Merton's model with $\lambda=3$ a year, what is the probability of at least one jump in one week (a year of 52
weeks)? In one year?
\end{exercise}

\begin{exercise}[$\star$]\label{exo:dv:jumps-and-levy-models:2}
Show that $\phi(-i)=1$ for the Merton characteristic function, and say what it means.
\end{exercise}

\begin{exercise}[$\star$]\label{exo:dv:jumps-and-levy-models:3}
A Lévy model has skewness $-0.4$ at one year. What is its skewness at one week and at four years?
\end{exercise}

\begin{exercise}[$\star\star$]\label{exo:dv:jumps-and-levy-models:4}
Derive the binomial hedge ratio across a two-outcome event, and compute it for the biotechnology straddle from the
values after the event: $V_u-V_0=2.21$, $V_d-V_0=-1.48$, moves $+22.1\%$ and $-14.8\%$ on a spot of 50.
\end{exercise}

\begin{exercise}[$\star\star$]\label{exo:dv:jumps-and-levy-models:5}
Why does the model's own delta of the straddle, $-0.02$, hedge worse than no hedge at all across the event?
\end{exercise}

\begin{exercise}[$\star\star$]\label{exo:dv:jumps-and-levy-models:6}
Explain why a Lévy model fitted to the one-month smile underestimates the one-year skew, and what Bates adds.
\end{exercise}

\begin{exercise}[$\star\star\star$]\label{exo:dv:jumps-and-levy-models:7}
\emph{Coding.} Re-run the daily hedging experiment with Merton's jumps halved in size and doubled in frequency
($\mu_J$ and $\delta$ halved, $\lambda$ doubled). How does the standard deviation of the hedged P\&L change?
\end{exercise}

\begin{exercise}[$\star\star\star$]\label{exo:dv:jumps-and-levy-models:8}
\emph{Find the flaw.} ``The one-week straddle trades at 160\% implied volatility. We sold it and delta-hedged with the
Black delta at 160\%, so our only risk is the difference between realised and implied volatility.''
\end{exercise}

\section{Problem: The Binary Event}

\begin{problem}[{Weekend problem --- what the one-week smile knows about Thursday}]\label{pb:dv:jumps-and-levy-models:1}
A share trades at 50. Its one-week smile, quoted at 13 strikes from 38 to 62, brackets the publication of trial
results. The desk's model is a diffusion plus one scheduled log-jump: up by $a$ with probability $p$, down by the
fair amount otherwise. The ``market'' is the chapter's richer model, in which the outcome sizes are themselves
uncertain.

\medskip\noindent\textbf{Part I --- The smile.}
\begin{enumerate}
\item Give the market's at-the-money implied volatility and its volatility at the 38 and 62 strikes.
\item Why does the smile fall on both sides of the money?
\item What \omterm{def:dv:parametrising-the-surface:event}{event variance} and standard deviation of the move does chapter 8's rule read from the at-the-money
volatility against a 35\% diffusion?
\item What expected absolute move does the rule give, and why is that the more reliable of its two numbers here?
\item Why can no single volatility price all 13 strikes?
\end{enumerate}

\medskip\noindent\textbf{Part II --- The fit.}
\begin{enumerate}[resume]
\item Fit $(\sigma,p,a)$: give the parameters and the error.
\item Give the implied up and down moves.
\item Why is the fitted diffusion volatility above the true 35\%?
\item What does the fitted model say the probability of success is, and how would a trader use that number?
\item What is the fair down move for $p=0.4$ and an up move of $+22.1\%$?
\end{enumerate}

\medskip\noindent\textbf{Part III --- The hedge.}
\begin{enumerate}[resume]
\item Give the straddle premium and the standard deviation of the unhedged P\&L across the event (known sizes).
\item Give the Black delta and the model delta of the straddle, and the P\&L standard deviation with each.
\item Give the binomial hedge ratio and the residual with known sizes.
\item Give the best ratio and the residual when the sizes are uncertain, in price and in percent of premium.
\item What instrument would hedge the rest?
\end{enumerate}

\medskip\noindent\textbf{Part IV --- Judgement.}
\begin{enumerate}[resume]
\item Why is the binomial ratio not the option's delta?
\item How should the desk hedge in the days before the event, and at the event?
\item What reserve would you hold on a short straddle across the event?
\item State the \emph{named result}: the jump size and probability implied by the one-week smile, and the error of
the delta hedge across the event.
\item In one sentence: what does a jump model give the desk that an implied volatility cannot?
\end{enumerate}
\end{problem}

\section{Interview questions}

\begin{interviewq}[$\star$ \iqroles{trader, researcher}]\label{iq:dv:jumps-and-levy-models:1}
Why do short-dated out-of-the-money options suggest jumps?
\end{interviewq}

\begin{interviewq}[$\star$ \iqroles{researcher}]\label{iq:dv:jumps-and-levy-models:2}
Write Merton's jump-diffusion and its price as a series of Black--Scholes prices.
\end{interviewq}

\begin{interviewq}[$\star\star$ \iqroles{trader, risk}]\label{iq:dv:jumps-and-levy-models:3}
Can you delta-hedge a jump? What is left, and how do you hedge it?
\end{interviewq}

\begin{interviewq}[$\star\star$ \iqroles{researcher}]\label{iq:dv:jumps-and-levy-models:4}
How do skewness and kurtosis scale with maturity in a Lévy model, and what does that imply for the smile?
\end{interviewq}

\begin{interviewq}[$\star\star$ \iqroles{developer}]\label{iq:dv:jumps-and-levy-models:5}
Your Fourier pricer is accurate for Heston at one year and wrong for variance gamma at one week. Why?
\end{interviewq}

\begin{interviewq}[$\star\star\star$ \iqroles{trader}]\label{iq:dv:jumps-and-levy-models:6}
A share reports earnings tomorrow. How do you price and hedge the one-week straddle, and what can go wrong?
\end{interviewq}
